\section*{Exercises about the Minimal Polynomial}

\exercise{1}
\begin{xrcs}
  Suppose $T \in \linmap(V)$. Prove that $9$ is an eigenvalue of $T^2$ $\iff$ $3$ or $-3$ is an eigenvalue of $T$.

  \begin{xprf}
    \Leftarrowdirection Suppose $-3$ is an eigenvalue of $T$. Thus, there exists $v \in V \setminus \{0\}$ such that $Tv = -3v$. Applying $T$ a second time yields $T^2 v = T(-3v) = -3 Tv = (-3)^2 v = 9v$. Therefore, $9$ is an eigenvalue of $T^2$. The case where $3$ is an eigenvalue of $T$ is identical.

    \Rightarrowdirection Let $v \in V \setminus \{0\}$ satisfy $T^2 v = 9v$. Then
    \begin{equation}
      (T^2 - 9I) v = (T - 3I)(T + 3I) v = 0.
    \end{equation}
    Set $w := (T + 3I)v$. Then $(T - 3I)w = 0$, so $Tw = 3w$.
    \begin{itemize}
      \item If $w \neq 0$, then $w$ is an eigenvector of $T$ corresponding to the eigenvalue $3$.
      \item If $w = 0$, then $(T + 3I)v = 0 \implies Tv = -3v$, so $v$ is an eigenvector of $T$ corresponding to the eigenvalue $-3$.
    \end{itemize}
    In either case, $3$ or $-3$ is an eigenvalue of $T$.
  \end{xprf}
\end{xrcs}

\exercise{3}
\begin{xrcs}
  Suppose $n \in \nat_{>0}$ and $T \in \linmap(\myF^n)$ is defined by $T(x_1 , \ddd, x_n) := (x_1 + \cdots + x_n, \ddd, x_1 + \cdots + x_n)$. The matrix $\mmatrix (T)$ of $T$ with respect to the standard basis is
  \begin{equation}
    \mmatrix (T) =
    \begin{pmatrix}
      1 & 1 & \cdots & 1 \\
      1 & 1 & \cdots & 1 \\
      \vdots & \vdots & \ddots & \vdots \\
      1 & 1 & \cdots & 1
    \end{pmatrix}.
  \end{equation}

  \begin{enumerate}
    \item Find all eigenvalues and eigenvectors of $T$.
    \item Find the minimal polynomial of $T$.
  \end{enumerate}

  \begin{xsol}
    We want to find all nonzero $x \in \myF^n$ such that $Tx = \lambda x$. The eigenvalue equation gives:
    \begin{equation}
      x_1 + \cdots + x_n = \lambda x_1 = \lambda x_2 = \cdots = \lambda x_n.
    \end{equation}
    Setting $S := x_1 + \cdots + x_n$, we have $S = \lambda x_k$ for each $k$. Summing over $k=1, \dots, n$ gives:
    \begin{equation}
      nS = \lambda S \iff S(n - \lambda) = 0.
    \end{equation}
    \begin{itemize}
      \item If $\lambda = n$, then $x_1 = x_2 = \cdots = x_n$. Thus, the eigenspace for $\lambda = n$ is the $1$-dimensional subspace spanned by $(1, 1, \ddd, 1)$.
      \item If $\lambda \neq n$, then $S = 0$, so $\lambda x_k = 0$ for all $k$, which requires $\lambda = 0$. The eigenspace for $\lambda = 0$ is the $(n-1)$-dimensional subspace consisting of all $(x_1, \dots, x_n)$ such that $x_1 + \cdots + x_n = 0$.
    \end{itemize}

    \prooffont{The minimal polynomial of $T$:}
    Let $v = (1, 1, \dots, 1)$. Since $\myrange T = \myspan{v}$ and $Tv = nv$, we have
    \begin{equation}
      T(T - nI)x = T(Tx - nx) = T(Sx_0 - nx) = 0 \quad \forall x \in \myF^n.
    \end{equation}
    Hence $T(T - nI) = 0$. Since $T \neq 0$ and $T \neq nI$ (for $n > 1$), the polynomial
    \begin{equation}
      p(z) = z(z - n)
    \end{equation}
    is the minimal polynomial of $T$. (If $n = 1$, $T = I$, so the minimal polynomial is $z - 1$).
  \end{xsol}
\end{xrcs}

\exercise{4}
\begin{xrcs}
  Suppose $\myF = \compl$, $T \in \linmap(V)$, $p \in \polyn(\compl)$, and $\alpha \in \compl$. Prove that $\alpha$ is an eigenvalue of $p(T)$ $\iff$ $\alpha = p(\lambda)$ for some eigenvalue $\lambda$ of $T$.

  \begin{xprf}
    \Leftarrowdirection Suppose $\lambda \in \compl$ is an eigenvalue of $T$. Then there exists a nonzero $v \in V$ such that $Tv = \lambda v$.
    By induction, $T^k v = \lambda^k v$ for every $k \in \nat$.
    For $p(z) = \sum_{j=0}^m c_j z^j$, we have
    \begin{equation}
      p(T)v = \sum_{j=0}^m c_j T^j v = \sum_{j=0}^m c_j \lambda^j v = p(\lambda) v.
    \end{equation}
    Since $v \neq 0$, $p(\lambda)$ is an eigenvalue of $p(T)$, so $\alpha = p(\lambda)$ is an eigenvalue.

    \Rightarrowdirection Suppose $\alpha$ is an eigenvalue of $p(T)$.
    Then the operator $p(T) - \alpha I$ is not injective.
    Since $\myF = \compl$, by the Fundamental Theorem of Algebra, the polynomial $p(z) - \alpha$ factors completely into linear factors over $\compl$:
    \begin{equation}
      p(z) - \alpha = c(z - \lambda_1)(z - \lambda_2)\cdots(z - \lambda_m)
    \end{equation}
    where $c \neq 0$ and $\lambda_1, \dots, \lambda_m \in \compl$.
    Substituting $T$ gives
    \begin{equation}
      p(T) - \alpha I = c(T - \lambda_1 I)(T - \lambda_2 I)\cdots(T - \lambda_m I).
    \end{equation}
    Since $p(T) - \alpha I$ is not injective, at least one of the factors $T - \lambda_j I$ must not be injective.
    Thus, $\lambda_j$ is an eigenvalue of $T$ for some $j \in \{1, \dots, m\}$.
    Evaluating $p(z) - \alpha$ at $z = \lambda_j$ gives
    \begin{equation}
      p(\lambda_j) - \alpha = 0 \implies \alpha = p(\lambda_j).
    \end{equation}
    This completes the proof.
  \end{xprf}
\end{xrcs}

% ==============================================================================
% ORIGINAL ATTEMPT (Commented out): Exercise 4 Draft
% ------------------------------------------------------------------------------
% AI FEEDBACK:
% 1. Logical Flaw: In the $\Rightarrow$ direction, the draft assumes an eigenvector $u$
%    of $p(T)$ satisfies $Tu = \lambda u$ (i.e., assuming $u$ is also an eigenvector of $T$).
%    This assumes the very claim to be proved.
% 2. Role of the Complex Field $\compl$: The draft asks "where have we been using the fact
%    that $\myF = \compl$?". Over $\compl$, every non-constant polynomial factors completely
%    into linear factors (Fundamental Theorem of Algebra). Factoring $p(z) - \alpha$ into
%    $c \prod (z - \lambda_j)$ gives $p(T) - \alpha I = c \prod (T - \lambda_j I)$, which fails
%    injectivity if and only if some factor $T - \lambda_j I$ fails injectivity.
% ==============================================================================
% \begin{xprf}
%   If $\lambda$ is an eigenvalue of $T$, then $\exists v \in V$ such that
%   \begin{align}
%     T v   &= \lambda   v \\
%     T^2 v &= \lambda^2 v \\
%     T^k v &= \lambda^k v \quad \forall k \in \nat_{>0}.
%   \end{align}
%
%   \Rightarrowdirection If $\alpha$ is an eigenvalue of $p(T)$, then $\exists u \in V$ such that
%   \begin{align}
%     \alpha u &= p(T)(u)  \\
%     \alpha u &= c_0 u + c_1 T u + \cdots + c_{m-1} T^{m-1} u + T^m u \\
%     \alpha u &= c_0 \lambda^0 u + c_1 \lambda^1 u + \cdots + c_{m-1} \lambda^{m-1} u + \lambda ^m u \\
%     \alpha u &= (c_0 \lambda^0 + c_1 \lambda^1 + \cdots + c_{m-1} \lambda^{m-1} + \lambda ^m ) u \\
%     \alpha   &= c_0 \lambda^0 + c_1 \lambda^1 + \cdots + c_{m-1} \lambda^{m-1} + \lambda ^m \\
%     \alpha   &= p(\lambda)
%   \end{align}
%
%   But where have we been using the fact that $\myF = \compl$?
% \end{xprf}

\exercise{7}
\begin{xrcs}
  Suppose $V$ is finite-dimensional and $T \in \linmap(V)$. Prove that $T$ is invertible if and only if the constant term of the minimal polynomial of $T$ is non-zero.
% \begin{xprf}
%   Let $p(z) = z^m + a_{m-1}z^{m-1} + \dots + a_1 z + a_0$ be the minimal polynomial of $T$.
%   By Theorem \ref{thm: eigenvalues are the zeros of the minimal polynomial}, $\lambda \in \myF$ is an eigenvalue of $T$ if and only if $p(\lambda) = 0$.
%   In particular, $0$ is an eigenvalue of $T \iff p(0) = 0 \iff a_0 = 0$.
%   Since $V$ is finite-dimensional, $T$ is invertible if and only if $0$ is not an eigenvalue of $T$ (i.e. $\mynull T = \{0\}$).
%   Therefore, $T$ is invertible if and only if $a_0 \neq 0$.
% \end{xprf}
\end{xrcs}

\exercise{12}
\begin{xrcs}
  Define $T \in \linmap(\myF^n)$ by $T(x_1, x_2, x_3, \ddd, x_n) = (x_1, 2 x_2, 3 x_3, \ddd,  n x_n)$. Find the minimal polynomial of $T$.

  \begin{xsol}
    For a standard basis $e_1, \ddd, e_n$ we have $T( e_1 + \cdots + e_n) = e_1 + 2e_2 + \cdots + n e_n$. Or equivalently:
    \begin{equation}
      \begin{aligned}
        T(e_1) &= e_1  \\
        T(e_2) &= 2e_2 \\
               &\; \vdots \\
        T(e_n) &= n e_n.
      \end{aligned}
    \end{equation}

    Hence, $e_1, \ddd, e_n$ are $\dim \myF^n = n$ distinct eigenvectors with corresponding eigenvalues $1, \ddd, n$. This makes $Te_1, \ddd, Te_n$ linearly independent as well. Using \ref{thm: eigenvalues are the zeros of the minimal polynomial}, we conclude that the minimal polynomial of $T$ looks like this:
    \begin{equation}
      \label{eq: minmal polynomial of T(x_1, x_2, x_3, ..., x_n) = (x_1, 2 x_2, 3 x_3, ...,  n x_n)}
      p(z) = (z-1)(z-2)\cdots(z-n) \quad \forall z \in \myF^n.
    \end{equation}

    We can check if $\forall v \in \myF^n$, given as $v = a_1 e_1 + \cdots + a_n e_n$ ($a_i \in \myF$), the following equation holds:
    \begin{equation}
      (T-I)(T-2I)\cdots(T-nI)v=0.
    \end{equation}

    This is true because $a_k e_k \in \mynull (T-kI)$ for $k=1, \ddd, n$. Therefore,
    \begin{equation}
      \begin{aligned}
          (T-I)(T-2I)\cdots(T-nI)v & = (T-I)(T-2I)(T-3I)\cdots(T-nI)(a_1 e_1) \\
                                 & \quad + (T-I)(T-2I)(T-3I)\cdots(T-nI)(a_2 e_2) \\
                                 & \quad + (T-I)(T-2I)(T-3I)\cdots(T-nI)(a_3 e_3) \\
                                 & \quad + \; \cdots \\
                                 & \quad + (T-I)(T-2I)(T-3I)\cdots(T-nI)(a_n e_n) = 0.
      \end{aligned}
    \end{equation}

    Thus, we have found our minimal polynomial given in equation \eqref{eq: minmal polynomial of T(x_1, x_2, x_3, ..., x_n) = (x_1, 2 x_2, 3 x_3, ...,  n x_n)}. It is easy to conclude that there exists no polynomial of smaller degree that distributes to every basis vector of the sum $a_1 e_1 + \cdots + a_n e_n$.
  \end{xsol}
\end{xrcs}

\exercise{16}
\begin{xrcs}
  Suppose $a_0, a_1, \ddd, a_n \in \myF$. Let $T$ be the operator on $\myF^n$ whose matrix (with respect to the standard basis) is
  \begin{equation}
    \mmatrix (T) =
    \begin{blockarray}{cccccccc}
               & e_1 & e_2 & \cdots & & e_{n-1} & e_n      \\
      \begin{block}{c(ccccccc)}
        e_1    & 0   &     &        & &         & -a_0     \\
        e_2    & 1   & 0   &        & &         & -a_1     \\
        e_3    &     & 1   & \ddots & &         & -a_2     \\
        \vdots &     &     & \ddots & &         &  \vdots  \\
        e_{n-1}&     &     &        & & 0       & -a_{n-2} \\
        e_n    &     &     &        & & 1       & -a_{n-1} \\
      \end{block}
    \end{blockarray}.
  \end{equation}

  Note that all the other entries of this matrix are zero. Show that the minimal polynomial of $T$ is the polynomial
  \begin{equation}
    a_0 + a_1 z + \cdots + a_{n-1} z^{n-1} + z^n.
  \end{equation}

  \begin{xsol}
    Looking at the matrix $\mmatrix(T)$ with respect to the standard basis $e_1, \dots, e_n$, we observe:
    \begin{equation}
      \begin{aligned}
        T e_1 &= e_2, \\
        T e_2 &= e_3 = T^2 e_1, \\
        T e_3 &= e_4 = T^3 e_1, \\
              & \; \; \vdots \\
        T e_{n-1} &= e_n = T^{n-1} e_1, \\
        T e_n &= -a_0 e_1 - a_1 e_2 - a_2 e_3 - \cdots - a_{n-1} e_n = T^n e_1.
      \end{aligned}
    \end{equation}
    Rearranging the last equation gives:
    \begin{equation}
      T^n e_1 + a_{n-1} T^{n-1} e_1 + \cdots + a_1 T e_1 + a_0 e_1 = 0 \iff (T^n + a_{n-1} T^{n-1} + \dots + a_0 I) e_1 = 0.
    \end{equation}
    The list $e_1, T e_1, \dots, T^{n-1} e_1 = e_1, e_2, \dots, e_n$ is linearly independent.
    Therefore, no nonzero polynomial of degree less than $n$ can annihilate $e_1$.
    Since $(T^n + a_{n-1}T^{n-1} + \dots + a_0 I) e_k = T^{k-1}(T^n + \dots + a_0 I)e_1 = 0$ for all $k \in \{1, \dots, n\}$, the polynomial
    \begin{equation}
      p(z) = z^n + a_{n-1} z^{n-1} + \cdots + a_1 z + a_0
    \end{equation}
    annihilates every basis vector and thus equals the minimal polynomial of $T$.
  \end{xsol}
\end{xrcs}

